% !TITLE 楚辞·九歌·山鬼
% !AUTHOR 屈原
% !DYNASTY 先秦
% !GENRE 楚辞
% !ORDER 9
% !SOURCE 楚辞·九歌

\begin{poem}
若有人兮山之阿\textsuperscript{1}，被薜荔兮带女萝\textsuperscript{2}。
既含睇兮又宜笑\textsuperscript{3}，子慕予兮善窈窕\textsuperscript{4}。
乘赤豹兮从文狸\textsuperscript{5}，辛夷车兮结桂旗\textsuperscript{6}。
被石兰兮带杜衡\textsuperscript{7}，折芳馨兮遗所思\textsuperscript{8}。
余处幽篁兮终不见天\textsuperscript{9}，路险难兮独后来\textsuperscript{10}。

表独立兮山之上\textsuperscript{11}，云容容兮而在下\textsuperscript{12}。
杳冥冥兮羌昼晦\textsuperscript{13}，东风飘兮神灵雨\textsuperscript{14}。
留灵修兮憺忘归\textsuperscript{15}，岁既晏兮孰华予\textsuperscript{16}。
采三秀兮于山间\textsuperscript{17}，石磊磊兮葛蔓蔓\textsuperscript{18}。
怨公子兮怅忘归\textsuperscript{19}，君思我兮不得闲\textsuperscript{20}。

山中人兮芳杜若\textsuperscript{21}，饮石泉兮荫松柏\textsuperscript{22}。
君思我兮然疑作\textsuperscript{23}。
雷填填兮雨冥冥\textsuperscript{24}，猿啾啾兮狖夜鸣\textsuperscript{25}。
风飒飒兮木萧萧\textsuperscript{26}，思公子兮徒离忧\textsuperscript{27}。
\end{poem}

\begin{annotations}
\notetext{1}{山之阿：山坳、山弯处。阿（ē），山弯曲的地方。好像有个人在山弯里若隐若现——山鬼出场，朦朦胧胧，不似凡人。}
\notetext{2}{被薜荔兮带女萝：披着薜荔做的衣裳，系着女萝做的腰带。被，同"披"；薜荔（bì lì），一种蔓生植物；女萝，松萝，寄生在树上的丝状植物。山鬼以草木为衣，写其野性与自然之美。}
\notetext{3}{含睇兮又宜笑：微微斜视，含着笑意。睇（dì），斜视、流盼；宜笑，笑得恰到好处。写山鬼眼神灵动、笑靥迷人。}
\notetext{4}{子慕予兮善窈窕：你爱慕我美好的身姿。子，你（指所思念的人）；慕，爱慕；窈窕（yǎo tiǎo），体态美好。}
\notetext{5}{乘赤豹兮从文狸：驾着赤豹拉的车，后面跟着花纹斑斓的狸猫。赤豹，红褐色的豹子；文狸，毛色有花纹的野猫。}
\notetext{6}{辛夷车兮结桂旗：辛夷木做的车，编结桂枝当旗帜。辛夷，即木兰，香木。}
\notetext{7}{被石兰兮带杜衡：披着石兰，系着杜衡。石兰，香草名；杜衡（héng），也是香草。山鬼周身都是芬芳的草木。}
\notetext{8}{折芳馨兮遗所思：折下芳香的花草，送给心中思念的人。遗（wèi），赠送。}
\notetext{9}{余处幽篁兮终不见天：我住在幽深的竹林里，始终看不见天。篁（huáng），竹林。}
\notetext{10}{路险难兮独后来：路太险太难了，所以我独自来迟了。山鬼迟到的原因——山路艰险，暗示人神之间的阻隔。}
\notetext{11}{表独立兮山之上：独自伫立在山顶之上。表，突出的、高高的。登高是为了看得更远，盼望那人到来。}
\notetext{12}{云容容兮而在下：云在脚下翻涌流淌。容容，云气浮动的样子。}
\notetext{13}{杳冥冥兮羌昼晦：深暗幽远，虽是白天却阴沉昏暗。杳（yǎo），深远；冥冥，昏暗；羌，乃、却；昼晦，白天变暗。}
\notetext{14}{东风飘兮神灵雨：东风吹来，神灵降下了雨。雨，这里用作动词——下雨。天气骤变，预示不祥。}
\notetext{15}{留灵修兮憺忘归：我想留住灵修（心上人），安然地忘了归去。灵修，楚辞中对所爱之人的美称（也暗喻君王）。憺（dàn），安然。}
\notetext{16}{岁既晏兮孰华予：年岁已经迟暮，谁还能让我重新容华焕发呢？晏，晚、暮；华，使……荣华。}
\notetext{17}{采三秀兮于山间：在山间采摘灵芝。三秀，灵芝草（一年开三次花，故称"三秀"）。灵芝是仙草，采之以待所思之人。}
\notetext{18}{石磊磊兮葛蔓蔓：山石堆积重叠，葛藤纠缠蔓延。磊磊（lěi），石块堆叠的样子；蔓蔓，藤蔓缠绕。山路难行的写照。}
\notetext{19}{怨公子兮怅忘归：怨恨公子啊，惆怅得忘了回去。}
\notetext{20}{君思我兮不得闲：你大概也想我吧，只是没有空闲。}
\notetext{21}{山中人兮芳杜若：我这山中之人像杜若一样芬芳。杜若，香草名。山鬼自比香草，重申自己的美好。}
\notetext{22}{饮石泉兮荫松柏：喝的是石缝中的泉水，遮荫的是松柏。写山鬼生活之高洁——饮清泉、居松柏之下。}
\notetext{23}{君思我兮然疑作：你思我？还是不信？信与疑在心中交替。然，相信；疑，怀疑；作，起。等待太久了，信心开始动摇。}
\notetext{24}{雷填填兮雨冥冥：雷声隆隆，雨色昏暗。填填，雷声。}
\notetext{25}{猿啾啾兮狖夜鸣：猿猴啾啾地啼叫，长尾猿在夜里悲鸣。狖（yòu），一种黑色长尾猿。雷雨交加，猿声哀鸣——烘托山鬼绝望的心情。}
\notetext{26}{风飒飒兮木萧萧：风声飒飒，树木萧萧。飒飒（sà），风声；萧萧，风吹树木声。}
\notetext{27}{思公子兮徒离忧：思念公子啊，只是白白地受着忧愁的折磨。徒，徒然、白白地；离，通"罹"，遭受。}
\end{annotations}

\begin{appreciation}
《山鬼》是《九歌》中最凄美的一首。表面上写的是山中女神等待恋人不至的故事，但那等待中的执着与渐次加深的绝望，超越了神与人的界限，成为所有等待者共同的心灵图景。但我选这首诗最主要的原因却很简单，这首诗里的“杜若”一词，就是你名字里“若”的来源。

山鬼出场的方式是朦胧的。"若有人"三个字，说明她不是一个确定的存在：似乎是个人，又似乎不是。她披着薜荔、系着女萝，驾赤豹、从文狸——周身都是山林的气息。她不是庙堂中端坐的庄严神祇，而是山野中自由生长的精灵。这种野性的美，是我们在《诗经》里所不曾见到的。

全诗的情感是渐次下沉的。开头时她满怀期待，折下芳草要送给心上人，迟到只是因为路太难走。她登上山顶，站在云海之上盼望。但天色变暗、风雨骤至，那人始终没有出现。于是她的心境从"憺忘归"（安然忘归）变成了"怅忘归"（惆怅忘归），从"君思我"的自我安慰变成了"君思我兮然疑作"的动摇。最后雷雨大作、猿声哀鸣，她只留下了一句"思公子兮徒离忧"——所有等待都是徒劳的，但即便如此，她还是在等。

屈原写《九歌》本是楚地的祭歌，但他的心灵投影无处不在。山鬼等待的"灵修"，在楚辞中常暗喻君王。屈原等待楚王的信任和召回，等了整整一生，从郢都等到了汨罗江——他终于没有等到。山鬼的"徒离忧"，正是屈原自己的命运。
\end{appreciation}
